\documentclass[10pt,letterpaper]{article}
\usepackage{cornell-notes}
\usepackage{mathtools}

\title{Chapter 4: Integration of Functions}
\author{}
\date{}

\setDocTitle{Chapter 4: Integration of Functions}
\setDocSubtitle{Numerical Methods}
\setDocVersion{v1.0}
\setDocStatus{Study Companion}
\setDocOwner{Numerical Methods Cornell Notes}
\setDocAuthor{}
\setDocDate{\today}
\setDocSummary{Cornell-format study and review notes for Chapter 4: Integration of Functions.}

\setCornellCollection{Numerical Methods Cornell Notes}
\setCornellUnitType{Chapter}
\setCornellUnitNumber{4}
\setCornellUnitTitle{Integration of Functions}
\setCornellCompanionLabel{Study and review companion}

\hypersetup{pdftitle={Chapter 4: Integration of Functions},pdfsubject={Cornell notes},pdfauthor={}}

\begin{document}
\maketitle

\section*{Chapter Overview}
This chapter organizes the mathematical and computational ideas behind integration of functions. The notes preserve the numbered section sequence while emphasizing method selection, explicit assumptions, numerical reliability, cost, and verification.

\section*{Learning Objectives}
Explain the governing problem class; compare Classical Formulas for Equally Spaced Abscissas, Elementary Algorithms, Romberg Integration, and Improper Integrals; design defensible stopping and verification checks; distinguish problem conditioning from algorithmic stability.

\section*{Prerequisites}
Calculus, Taylor expansions, interpolation, and error tolerances.

\section*{Operational Workflow}
Analyze smoothness and domain; transform singular or infinite regions; choose fixed, Gaussian, adaptive, or stochastic quadrature; allocate error and evaluation budgets; verify by refinement.

\subsection*{Formula and notation anchor}
\[
\int_a^b f(x)\,dx\approx\sum_{i=1}^{m} w_i f(x_i)
\]
Nodes and weights encode the rule; the error depends on which function class the rule integrates exactly and on unresolved behavior.

\begin{warnbox}{Numerical warning}
Endpoint singularities, cancellation, missed narrow features, false convergence, excessive recursion, and the curse of dimensionality.
\end{warnbox}

\section{4.0 Introduction}
\begin{cornell}
\noterow{What is the central idea?}{Numerical quadrature replaces an integral by weighted samples; accuracy depends on smoothness, interval geometry, singularities, and error control.}
\noterow{How should it be used?}{For Introduction, begin from explicit inputs, outputs, preconditions, and representation choices.  Analyze smoothness and domain; transform singular or infinite regions; choose fixed, Gaussian, adaptive, or stochastic quadrature; allocate error and evaluation budgets; verify by refinement.}
\noterow{Cost and reliability}{Analyze arithmetic work, storage, data movement, and convergence separately.  Relevant failure pressures include endpoint singularities, cancellation, missed narrow features, false convergence, excessive recursion, and the curse of dimensionality.  Use scaled tolerances and report both success criteria and diagnostic evidence.}
\noterow{Implementation checklist}{\begin{itemize}
\item Validate dimensions, domains, ordering, and structural assumptions.
\item Guard small divisors, invalid states, overflow, underflow, and nonfinite values.
\item Make mutation, workspace, indexing, and termination behavior explicit.
\item Test a simple exact case, a difficult valid case, a boundary case, and an expected failure.
\end{itemize}}
\end{cornell}
\begin{summarybox}{Section summary}
Numerical quadrature replaces an integral by weighted samples; accuracy depends on smoothness, interval geometry, singularities, and error control.  Select this technique only when its assumptions match the data, and accept a result only after an independent residual, error estimate, invariant, or refinement check.
\end{summarybox}

\section{4.1 Classical Formulas for Equally Spaced Abscissas}
\begin{cornell}
\noterow{What is the central idea?}{Newton-Cotes rules integrate an interpolating polynomial on equally spaced nodes, with truncation error tied to higher derivatives.}
\noterow{How should it be used?}{For Classical Formulas for Equally Spaced Abscissas, begin from explicit inputs, outputs, preconditions, and representation choices.  Analyze smoothness and domain; transform singular or infinite regions; choose fixed, Gaussian, adaptive, or stochastic quadrature; allocate error and evaluation budgets; verify by refinement.}
\noterow{Quantitative lens}{Identify the governing equation, recurrence, probability law, or geometric predicate for Classical Formulas for Equally Spaced Abscissas.  Define every variable and scale; then derive a residual, error estimate, or invariant that can be checked independently.}
\noterow{Cost and reliability}{Analyze arithmetic work, storage, data movement, and convergence separately.  Relevant failure pressures include endpoint singularities, cancellation, missed narrow features, false convergence, excessive recursion, and the curse of dimensionality.  Use scaled tolerances and report both success criteria and diagnostic evidence.}
\noterow{Implementation checklist}{\begin{itemize}
\item Validate dimensions, domains, ordering, and structural assumptions.
\item Guard small divisors, invalid states, overflow, underflow, and nonfinite values.
\item Make mutation, workspace, indexing, and termination behavior explicit.
\item Test a simple exact case, a difficult valid case, a boundary case, and an expected failure.
\end{itemize}}
\end{cornell}
\begin{summarybox}{Section summary}
Newton-Cotes rules integrate an interpolating polynomial on equally spaced nodes, with truncation error tied to higher derivatives.  Select this technique only when its assumptions match the data, and accept a result only after an independent residual, error estimate, invariant, or refinement check.
\end{summarybox}

\section{4.2 Elementary Algorithms}
\begin{cornell}
\noterow{What is the central idea?}{Composite trapezoid and Simpson procedures refine a partition and reuse evaluations, creating practical building blocks for adaptive integration.}
\noterow{How should it be used?}{For Elementary Algorithms, begin from explicit inputs, outputs, preconditions, and representation choices.  Analyze smoothness and domain; transform singular or infinite regions; choose fixed, Gaussian, adaptive, or stochastic quadrature; allocate error and evaluation budgets; verify by refinement.}
\noterow{Quantitative lens}{Identify the governing equation, recurrence, probability law, or geometric predicate for Elementary Algorithms.  Define every variable and scale; then derive a residual, error estimate, or invariant that can be checked independently.}
\noterow{Cost and reliability}{Analyze arithmetic work, storage, data movement, and convergence separately.  Relevant failure pressures include endpoint singularities, cancellation, missed narrow features, false convergence, excessive recursion, and the curse of dimensionality.  Use scaled tolerances and report both success criteria and diagnostic evidence.}
\noterow{Implementation checklist}{\begin{itemize}
\item Validate dimensions, domains, ordering, and structural assumptions.
\item Guard small divisors, invalid states, overflow, underflow, and nonfinite values.
\item Make mutation, workspace, indexing, and termination behavior explicit.
\item Test a simple exact case, a difficult valid case, a boundary case, and an expected failure.
\end{itemize}}
\end{cornell}
\begin{summarybox}{Section summary}
Composite trapezoid and Simpson procedures refine a partition and reuse evaluations, creating practical building blocks for adaptive integration.  Select this technique only when its assumptions match the data, and accept a result only after an independent residual, error estimate, invariant, or refinement check.
\end{summarybox}

\section{4.3 Romberg Integration}
\begin{cornell}
\noterow{What is the central idea?}{Romberg integration extrapolates a sequence of trapezoidal estimates to cancel leading even-power error terms.}
\noterow{How should it be used?}{For Romberg Integration, begin from explicit inputs, outputs, preconditions, and representation choices.  Analyze smoothness and domain; transform singular or infinite regions; choose fixed, Gaussian, adaptive, or stochastic quadrature; allocate error and evaluation budgets; verify by refinement.}
\noterow{Quantitative anchor}{\[
R_{k,j}=R_{k,j-1}+\frac{R_{k,j-1}-R_{k-1,j-1}}{4^j-1}
\]
State the dimensions, scaling, and validity assumptions before using this relation.  Check the resulting residual or invariant rather than trusting an iteration count alone.}
\noterow{Cost and reliability}{Analyze arithmetic work, storage, data movement, and convergence separately.  Relevant failure pressures include endpoint singularities, cancellation, missed narrow features, false convergence, excessive recursion, and the curse of dimensionality.  Use scaled tolerances and report both success criteria and diagnostic evidence.}
\noterow{Implementation checklist}{\begin{itemize}
\item Validate dimensions, domains, ordering, and structural assumptions.
\item Guard small divisors, invalid states, overflow, underflow, and nonfinite values.
\item Make mutation, workspace, indexing, and termination behavior explicit.
\item Test a simple exact case, a difficult valid case, a boundary case, and an expected failure.
\end{itemize}}
\end{cornell}
\begin{summarybox}{Section summary}
Romberg integration extrapolates a sequence of trapezoidal estimates to cancel leading even-power error terms.  Select this technique only when its assumptions match the data, and accept a result only after an independent residual, error estimate, invariant, or refinement check.
\end{summarybox}

\section{4.4 Improper Integrals}
\begin{cornell}
\noterow{What is the central idea?}{Infinite intervals and endpoint singularities should be transformed or partitioned so that the numerical rule sees a finite, sufficiently regular integrand.}
\noterow{How should it be used?}{For Improper Integrals, begin from explicit inputs, outputs, preconditions, and representation choices.  Analyze smoothness and domain; transform singular or infinite regions; choose fixed, Gaussian, adaptive, or stochastic quadrature; allocate error and evaluation budgets; verify by refinement.}
\noterow{Quantitative lens}{Identify the governing equation, recurrence, probability law, or geometric predicate for Improper Integrals.  Define every variable and scale; then derive a residual, error estimate, or invariant that can be checked independently.}
\noterow{Cost and reliability}{Analyze arithmetic work, storage, data movement, and convergence separately.  Relevant failure pressures include endpoint singularities, cancellation, missed narrow features, false convergence, excessive recursion, and the curse of dimensionality.  Use scaled tolerances and report both success criteria and diagnostic evidence.}
\noterow{Implementation checklist}{\begin{itemize}
\item Validate dimensions, domains, ordering, and structural assumptions.
\item Guard small divisors, invalid states, overflow, underflow, and nonfinite values.
\item Make mutation, workspace, indexing, and termination behavior explicit.
\item Test a simple exact case, a difficult valid case, a boundary case, and an expected failure.
\end{itemize}}
\end{cornell}
\begin{summarybox}{Section summary}
Infinite intervals and endpoint singularities should be transformed or partitioned so that the numerical rule sees a finite, sufficiently regular integrand.  Select this technique only when its assumptions match the data, and accept a result only after an independent residual, error estimate, invariant, or refinement check.
\end{summarybox}

\section{4.5 Quadrature by Variable Transformation}
\begin{cornell}
\noterow{What is the central idea?}{A change of variables redistributes sample density and can regularize singular or infinite-domain integrals before applying standard quadrature.}
\noterow{How should it be used?}{For Quadrature by Variable Transformation, begin from explicit inputs, outputs, preconditions, and representation choices.  Analyze smoothness and domain; transform singular or infinite regions; choose fixed, Gaussian, adaptive, or stochastic quadrature; allocate error and evaluation budgets; verify by refinement.}
\noterow{Quantitative lens}{Identify the governing equation, recurrence, probability law, or geometric predicate for Quadrature by Variable Transformation.  Define every variable and scale; then derive a residual, error estimate, or invariant that can be checked independently.}
\noterow{Cost and reliability}{Analyze arithmetic work, storage, data movement, and convergence separately.  Relevant failure pressures include endpoint singularities, cancellation, missed narrow features, false convergence, excessive recursion, and the curse of dimensionality.  Use scaled tolerances and report both success criteria and diagnostic evidence.}
\noterow{Implementation checklist}{\begin{itemize}
\item Validate dimensions, domains, ordering, and structural assumptions.
\item Guard small divisors, invalid states, overflow, underflow, and nonfinite values.
\item Make mutation, workspace, indexing, and termination behavior explicit.
\item Test a simple exact case, a difficult valid case, a boundary case, and an expected failure.
\end{itemize}}
\end{cornell}
\begin{summarybox}{Section summary}
A change of variables redistributes sample density and can regularize singular or infinite-domain integrals before applying standard quadrature.  Select this technique only when its assumptions match the data, and accept a result only after an independent residual, error estimate, invariant, or refinement check.
\end{summarybox}

\section{4.6 Gaussian Quadratures and Orthogonal Polynomials}
\begin{cornell}
\noterow{What is the central idea?}{Gaussian quadrature selects nodes and weights from orthogonal polynomials to integrate the highest possible polynomial degree for a fixed sample count.}
\noterow{How should it be used?}{For Gaussian Quadratures and Orthogonal Polynomials, begin from explicit inputs, outputs, preconditions, and representation choices.  Analyze smoothness and domain; transform singular or infinite regions; choose fixed, Gaussian, adaptive, or stochastic quadrature; allocate error and evaluation budgets; verify by refinement.}
\noterow{Quantitative lens}{Identify the governing equation, recurrence, probability law, or geometric predicate for Gaussian Quadratures and Orthogonal Polynomials.  Define every variable and scale; then derive a residual, error estimate, or invariant that can be checked independently.}
\noterow{Cost and reliability}{Analyze arithmetic work, storage, data movement, and convergence separately.  Relevant failure pressures include endpoint singularities, cancellation, missed narrow features, false convergence, excessive recursion, and the curse of dimensionality.  Use scaled tolerances and report both success criteria and diagnostic evidence.}
\noterow{Implementation checklist}{\begin{itemize}
\item Validate dimensions, domains, ordering, and structural assumptions.
\item Guard small divisors, invalid states, overflow, underflow, and nonfinite values.
\item Make mutation, workspace, indexing, and termination behavior explicit.
\item Test a simple exact case, a difficult valid case, a boundary case, and an expected failure.
\end{itemize}}
\end{cornell}
\begin{summarybox}{Section summary}
Gaussian quadrature selects nodes and weights from orthogonal polynomials to integrate the highest possible polynomial degree for a fixed sample count.  Select this technique only when its assumptions match the data, and accept a result only after an independent residual, error estimate, invariant, or refinement check.
\end{summarybox}

\section{4.7 Adaptive Quadrature}
\begin{cornell}
\noterow{What is the central idea?}{Adaptive quadrature estimates local error, subdivides difficult regions, and stops when a global tolerance budget is met.}
\noterow{How should it be used?}{For Adaptive Quadrature, begin from explicit inputs, outputs, preconditions, and representation choices.  Analyze smoothness and domain; transform singular or infinite regions; choose fixed, Gaussian, adaptive, or stochastic quadrature; allocate error and evaluation budgets; verify by refinement.}
\noterow{Quantitative anchor}{\[
|I_{\mathrm{fine}}-I_{\mathrm{coarse}}|\le \tau_{\mathrm{local}}
\]
State the dimensions, scaling, and validity assumptions before using this relation.  Check the resulting residual or invariant rather than trusting an iteration count alone.}
\noterow{Cost and reliability}{Analyze arithmetic work, storage, data movement, and convergence separately.  Relevant failure pressures include endpoint singularities, cancellation, missed narrow features, false convergence, excessive recursion, and the curse of dimensionality.  Use scaled tolerances and report both success criteria and diagnostic evidence.}
\noterow{Implementation checklist}{\begin{itemize}
\item Validate dimensions, domains, ordering, and structural assumptions.
\item Guard small divisors, invalid states, overflow, underflow, and nonfinite values.
\item Make mutation, workspace, indexing, and termination behavior explicit.
\item Test a simple exact case, a difficult valid case, a boundary case, and an expected failure.
\end{itemize}}
\end{cornell}
\begin{summarybox}{Section summary}
Adaptive quadrature estimates local error, subdivides difficult regions, and stops when a global tolerance budget is met.  Select this technique only when its assumptions match the data, and accept a result only after an independent residual, error estimate, invariant, or refinement check.
\end{summarybox}

\section{4.8 Multidimensional Integrals}
\begin{cornell}
\noterow{What is the central idea?}{Multidimensional integration must balance tensor-product cost, adaptive partitioning, and stochastic alternatives as dimension increases.}
\noterow{How should it be used?}{For Multidimensional Integrals, begin from explicit inputs, outputs, preconditions, and representation choices.  Analyze smoothness and domain; transform singular or infinite regions; choose fixed, Gaussian, adaptive, or stochastic quadrature; allocate error and evaluation budgets; verify by refinement.}
\noterow{Quantitative lens}{Identify the governing equation, recurrence, probability law, or geometric predicate for Multidimensional Integrals.  Define every variable and scale; then derive a residual, error estimate, or invariant that can be checked independently.}
\noterow{Cost and reliability}{Analyze arithmetic work, storage, data movement, and convergence separately.  Relevant failure pressures include endpoint singularities, cancellation, missed narrow features, false convergence, excessive recursion, and the curse of dimensionality.  Use scaled tolerances and report both success criteria and diagnostic evidence.}
\noterow{Implementation checklist}{\begin{itemize}
\item Validate dimensions, domains, ordering, and structural assumptions.
\item Guard small divisors, invalid states, overflow, underflow, and nonfinite values.
\item Make mutation, workspace, indexing, and termination behavior explicit.
\item Test a simple exact case, a difficult valid case, a boundary case, and an expected failure.
\end{itemize}}
\end{cornell}
\begin{summarybox}{Section summary}
Multidimensional integration must balance tensor-product cost, adaptive partitioning, and stochastic alternatives as dimension increases.  Select this technique only when its assumptions match the data, and accept a result only after an independent residual, error estimate, invariant, or refinement check.
\end{summarybox}

\section*{Method comparison}
\begin{center}
\small
\begin{tabularx}{\linewidth}{>{\bfseries\raggedright\arraybackslash}p{0.22\linewidth}XX}
\toprule
Method or lens & Best use & Main caution \\
\midrule
Composite Simpson & Smooth one-dimensional data & Fourth-order globally under standard assumptions \\
Gaussian quadrature & Smooth weighted integrands & High polynomial exactness; fixed nodes \\
Adaptive quadrature & Localized difficulty & Needs a trustworthy local error estimate \\
\bottomrule
\end{tabularx}
\end{center}

\section*{Worked example}
\begin{exambox}{Independent worked example}
Approximate $\int_0^1 x^2\,dx$ with Simpson's rule: $\tfrac16[f(0)+4f(1/2)+f(1)]=\tfrac16(0+1+1)=1/3$.  The result matches the exact integral because the rule is exact for cubics.  This success does not transfer to a singular integrand unless its regularity assumptions are restored.
\end{exambox}

\section*{Chapter synthesis}
\begin{summarybox}{Chapter synthesis}
The chapter's methods should be viewed as a decision system rather than a list of routines.  Begin with the mathematical problem and its conditioning, exploit structure, select an algorithm with compatible stability and cost, and verify the computed answer with evidence that is independent of the internal iteration.  The recurring implementation discipline is: analyze smoothness and domain; transform singular or infinite regions; choose fixed, gaussian, adaptive, or stochastic quadrature; allocate error and evaluation budgets; verify by refinement.
\end{summarybox}

\subsection*{Common mistakes and numerical hazards}
\begin{warnbox}{Common mistakes and numerical hazards}
Endpoint singularities, cancellation, missed narrow features, false convergence, excessive recursion, and the curse of dimensionality.  A second recurring mistake is reporting only a computed value: always retain tolerances, iteration counts, residuals, refinement behavior, and relevant structural checks.
\end{warnbox}

\subsection*{Retrieval practice}
\textbf{Questions}
\begin{enumerate}
\item What problem does Classical Formulas for Equally Spaced Abscissas address, and which assumption is easiest to violate?
\item What problem does Elementary Algorithms address, and which assumption is easiest to violate?
\item What problem does Romberg Integration address, and which assumption is easiest to violate?
\item What problem does Adaptive Quadrature address, and which assumption is easiest to violate?
\item What problem does Multidimensional Integrals address, and which assumption is easiest to violate?
\item How do conditioning, stability, and the final residual play different roles in this chapter?
\end{enumerate}

\textbf{Brief answers}
\begin{enumerate}
\item Newton-Cotes rules integrate an interpolating polynomial on equally spaced nodes, with truncation error tied to higher derivatives. Recheck the section's domain, structure, scaling, and failure diagnostics.
\item Composite trapezoid and Simpson procedures refine a partition and reuse evaluations, creating practical building blocks for adaptive integration. Recheck the section's domain, structure, scaling, and failure diagnostics.
\item Romberg integration extrapolates a sequence of trapezoidal estimates to cancel leading even-power error terms. Recheck the section's domain, structure, scaling, and failure diagnostics.
\item Adaptive quadrature estimates local error, subdivides difficult regions, and stops when a global tolerance budget is met. Recheck the section's domain, structure, scaling, and failure diagnostics.
\item Multidimensional integration must balance tensor-product cost, adaptive partitioning, and stochastic alternatives as dimension increases. Recheck the section's domain, structure, scaling, and failure diagnostics.
\item Conditioning describes sensitivity of the mathematical problem, stability describes error propagation by the algorithm, and the residual measures satisfaction of the computed equations; none is a universal substitute for the others.
\end{enumerate}

\subsection*{Mastery checklist}
\begin{itemize}
\item[$\square$] I can explain and select Classical Formulas for Equally Spaced Abscissas without relying on a memorized code listing.
\item[$\square$] I can explain and select Elementary Algorithms without relying on a memorized code listing.
\item[$\square$] I can explain and select Romberg Integration without relying on a memorized code listing.
\item[$\square$] I can explain and select Improper Integrals without relying on a memorized code listing.
\item[$\square$] I can explain and select Quadrature by Variable Transformation without relying on a memorized code listing.
\item[$\square$] I can state a scaled stopping rule and an independent verification check.
\item[$\square$] I can identify at least one conditioning risk and one algorithmic stability risk.
\item[$\square$] I can estimate time, storage, and data-movement costs at the appropriate level.
\end{itemize}

\end{document}

